\chapter{Operator kernels and traces}
\label{chap:operator-kernels}

The identity on an infinite-dimensional field space cannot be a finite sum
of rank-one operators. A completed kernel can represent it, but this does
not ensure that its trace exists. For formal Hamiltonian fields, these
questions have different answers. Evaluation and composition extend to
the completed kernels. The finite-rank trace does not extend continuously,
and no nonzero cyclic trace exists on the full operator algebra.

The adjoint action gives a concrete infinite-dimensional matter candidate.
Its operator kernels must retain all formal Hamiltonians, including
translations. The first task is therefore to complete the kernels without
replacing the symmetry algebra by finite Lie quotients.

\section{A kernel for each finite output}

Fix a discrete field $k$ of characteristic zero. Let
\[
 H=k[[x,y]]/k=\prod_{d\geq1}H_d,
 \qquad U_N=\prod_{d>N}H_d,\qquad H_N=H/U_N,
\]
where $H_d$ is the space of homogeneous polynomials of degree $d$.
The topology on $H$ has the $U_N$ as neighborhoods of zero. These
quotients are finite-dimensional vector spaces. Their use here does not
require quotient Lie brackets.

Index the nonconstant monomials by a countable set $\mathcal I$,
writing $e_I=x^ay^b$ when $I=(a,b)$ and $a+b\geq1$.
Their coefficient functionals form the continuous dual
\[
 D=\operatorname{Hom}_{\mathrm{cont}}(H,k)
   =\bigoplus_{I\in\mathcal I}k\epsilon^I.
\]
Give $D$ the discrete topology. A functional is continuous precisely
when it uses finitely many coefficients: its kernel must contain some
$U_N$, and this condition is also sufficient.

An elementary tensor $v\otimes\lambda$ acts by
$h\mapsto\lambda(h)v$. A finite sum of these tensors has finite-dimensional
image. To represent more operators, allow the finite set of input
functionals to depend on the output coefficients being retained.

Let $V$ be a complete Hausdorff vector space with a decreasing countable
basis of open linear subspaces $V_N^0$. Assume that
$V_N=V/V_N^0$ is finite-dimensional. Completeness here means that
compatible elements of all the quotients determine a unique element
of $V$. To form coefficients on $H$, introduce a commuting symbol $u_I$
for each monomial $e_I$. Both $H$ and the coefficient algebra
\[
 A=\prod_\alpha ku^\alpha
\]
have this form. In $A$, the index $\alpha$ ranges over finitely supported
families of nonnegative integers on $\mathcal I$. One can retain all
monomials of polynomial degree at most $N$ whose labels have degree
at most $N$. These finite sets exhaust the coefficient indices.

Give $V\otimes_kD$ the topology with neighborhoods $V_N^0\otimes_kD$.
Define
\begin{equation}
 K(V,D)=\varprojlim_N(V_N\otimes_kD)
       =\varprojlim_N\varinjlim_M(V_N\otimes_kH_M^\vee).
 \label{eq:op-completion}
\end{equation}
Each $V_N\otimes D$ is discrete. It need not be finite-dimensional.

\begin{proposition}
The space $K(V,D)$ is the complete Hausdorff completion of the stated
topology on $V\otimes D$. The algebraic tensor product embeds densely.
\end{proposition}
\begin{proof}
Write a tensor as $t=\sum_{j=1}^r v_j\otimes\lambda_j$ with the
$\lambda_j$ linearly independent. If every quotient image of $t$
vanishes, linear independence forces $v_j\in V_N^0$ for every $N$.
Separation gives $v_j=0$. This proves injectivity.

The quotient map onto $V_N\otimes D$ is surjective and has kernel
$V_N^0\otimes D$. Every neighborhood in the inverse limit can be
specified at one sufficiently large $N$. Surjectivity therefore proves
density and identifies the induced topology. Finally, compatible
sequences form a closed subspace of the product of discrete complete
spaces $V_N\otimes D$. Their inverse limit is complete and Hausdorff.
\end{proof}

\section{Continuous operators and the order of limits}

Write $\mathcal L(H,V)$ for continuous $k$-linear maps $H\to V$.
For the quotient map $q_N:V\to V_N$, give this space neighborhoods
\[
 \mathcal N_N=\{T:q_NT=0\}
             =\{T:T(H)\subset V_N^0\}.
\]
Thus two nearby operators agree on their entire output in a prescribed
finite quotient, for every input.

\begin{proposition}
Evaluation of tensors induces a topological isomorphism
\begin{equation}
 K(V,D)\xrightarrow{\ \sim\ }\mathcal L(H,V).
 \label{eq:op-hom}
\end{equation}
\end{proposition}
\begin{proof}
Choose a basis $v_1,\ldots,v_r$ of $V_N$. Every continuous map to
this discrete finite-dimensional space has the form
\[
 T_N(h)=\sum_{j=1}^r\lambda_j(h)v_j,
 \qquad \lambda_j\in D.
\]
Conversely, this formula defines a continuous map. Uniqueness of its
coordinate functionals gives
$V_N\otimes D=\operatorname{Hom}_{\mathrm{cont}}(H,V_N)$.
Compatible maps $T_N$ determine $T:H\to V$ by completeness of $V$.
The map $T$ is continuous because every $q_NT$ is continuous.
Its quotient kernels are exactly the prescribed operator neighborhoods.
\end{proof}

Enumerate the monomials in increasing total degree as $e_0,e_1,\ldots$.
Then $H\cong k^{\mathbb N}$, and an endomorphism has a matrix
\begin{equation}
 (Th)_i=\sum_jT_{ij}h_j,
 \qquad \#\{j:T_{ij}\ne0\}<\infty\quad\text{for each }i.
 \label{eq:op-rows}
\end{equation}
Such a matrix is called row-finite. Each output coefficient is a
continuous functional, so this condition is necessary and sufficient.
The operator topology makes each entire row discrete. It permits
different finite input supports in different rows.

As an identity of vector spaces,
\[
 V\otimes D
   =\varinjlim_M\varprojlim_N(V_N\otimes H_M^\vee).
\]
These maps factor through one common input quotient $H_M$.
Reversing the limits gives~\eqref{eq:op-completion}, where the required
$M$ can increase with $N$. The displayed equality concerns vector spaces.
One can instead give this union the final linear topology: a linear
subspace is open when its intersection with each finite-$M$ stage is open.
The trace below distinguishes that topology from the tensor topology.

The identity operator belongs to the latter completion:
\begin{equation}
 \mathbf1_H=\sum_{i\geq0}e_i\otimes\epsilon^i.
 \label{eq:op-identity}
\end{equation}
Its image in $H_N\otimes D$ is a finite sum. These images are
compatible, which is exactly what the completed expression means.

\section{Which bilinear contractions extend}

The topology used for kernels is stronger than the condition that a
bilinear map be jointly continuous. This distinction decides whether
closing an operator kernel to a scalar is possible.

\begin{proposition}
Let $W$ be a complete Hausdorff linearly topologized vector space,
meaning that its topology has a basis of open linear subspaces.
Continuous linear maps $K(V,D)\to W$ correspond to bilinear maps
$b:V\times D\to W$ satisfying
\begin{equation}
 \text{for every open linear }W^0\subset W,\quad
 \text{some }N\text{ has }b(V_N^0,D)\subset W^0.
 \label{eq:op-uniform}
\end{equation}
\end{proposition}
\begin{proof}
The linearization of $b$ on $V\otimes D$ is continuous exactly when
\eqref{eq:op-uniform} holds. Modulo $W^0$, it then factors through
$V_N\otimes D$. Apply this quotient map to the $N$-th component of
an element of $K(V,D)$. Increasing $N$ gives the same answer.
The answers for different $W^0$ are compatible. Completeness and
separation of $W$ give a unique element of $W$. This construction
defines a continuous extension. Density gives uniqueness. Conversely,
the inverse image of $W^0$ under a continuous linear extension contains
one of the defining kernel neighborhoods, giving~\eqref{eq:op-uniform}.
\end{proof}

Condition~\eqref{eq:op-uniform} is uniform over the entire second
argument. It implies joint continuity: near $(v,\lambda)$, keep
$\lambda$ fixed using discreteness of $D$, and vary $v$ in $V_N^0$.
Joint continuity alone does not imply this condition. For example,
\[
 H\times D\longrightarrow k,\qquad(h,\lambda)\longmapsto\lambda(h)
\]
is jointly continuous. For each fixed $\lambda$, its kernel in $H$
is open. Uniformity would instead require an open $U\subset H$
annihilated by every element of $D$. Since the coefficient functionals
separate points, $U$ would be zero. The zero subspace is not open in $H$.

\section{Evaluation, composition, and Hamiltonian symmetry}

Evaluation $\mathcal L(H,V)\times H\to V$ is jointly continuous.
Fix $(T_0,h_0)$ and an open output subspace $V^0$. Require
\[
 \Delta T(H)\subset V^0,\qquad
 \Delta h\in T_0^{-1}(V^0).
\]
Then
\[
 (T_0+\Delta T)(h_0+\Delta h)-T_0h_0
   =T_0\Delta h+\Delta T(h_0+\Delta h)\in V^0.
\]

Composition is jointly continuous as well. Give each continuous Hom
space the topology of agreement in output quotients. For
$S_0:V\to W$, choose an open $V^0$ with $S_0(V^0)\subset W^0$.
The conditions
\[
 \Delta S(V)\subset W^0,\qquad\Delta T(H)\subset V^0
\]
make every term of
\[
 (S_0+\Delta S)(T_0+\Delta T)-S_0T_0
 =S_0\Delta T+\Delta S\,T_0+\Delta S\,\Delta T
\]
map $H$ into $W^0$. In particular,
$K(H,D)=\operatorname{End}_{\mathrm{cont}}(H)$ is a complete
topological associative algebra. On elementary kernels,
\[
 (v\otimes\lambda)(w\otimes\mu)
       =\lambda(w)v\otimes\mu.
\]

Keeping whole rows fixed is necessary for this continuity statement.
In the weaker topology of entrywise matrix convergence, both
\[
 A_n=e_0\otimes\epsilon^n,\qquad
 B_n=e_n\otimes\epsilon^0
\]
tend to zero, but $A_nB_n=e_0\otimes\epsilon^0$.
In the operator topology used here, $A_n$ does not tend to zero.

The Hamiltonian bracket on $H$ is $[f,g]=f_xg_y-f_yg_x$ modulo constants.
Let $a_f=\operatorname{ad}_f$. The map
$f\mapsto a_f$ is continuous into $\operatorname{End}_{\mathrm{cont}}(H)$:
an entire bracket output through degree $N$ uses $f$ only through
degree $N+1$, uniformly over the other input.

For $V=A$, let $\rho_A(f)$ be the coefficient action
$\rho_A(f)O_g=O_{[f,g]}$, where $O_g=\sum_Ig_Iu_I$.
It is also continuous as a map into the operator space. For a fixed
output monomial $u^\alpha$, the action replaces a variable $u_J$
by $u_K$, with $K$ in the finite support of $\alpha$. Its coefficient
requires $|I|+|J|=|K|+2$. Only finitely many $I,J,K$ and predecessor
monomials $u^{\alpha-\mathbf e_K+\mathbf e_J}$ can occur.
Thus each entire output row depends on finitely many coefficients of $f$.

For $V=H$, put $\rho_V(f)=a_f$. The diagonal tensor action is
\[
 f\cdot(v\otimes\lambda)
  =\rho_V(f)v\otimes\lambda+v\otimes(f\cdot\lambda),
 \qquad(f\cdot\lambda)(g)=-\lambda([f,g]).
\]
Under~\eqref{eq:op-hom}, this action becomes
\begin{equation}
 f\cdot T=\rho_V(f)T-Ta_f.
 \label{eq:op-diagonal}
\end{equation}
The formula follows by evaluating an elementary tensor. It extends
jointly continuously by the proved composition and representation
properties. Expanding two actions and using Jacobi gives
$[f,g]\cdot T=f\cdot(g\cdot T)-g\cdot(f\cdot T)$.
No finite output quotient needs to carry a Hamiltonian action.

\section{Derivations and the Euler kernel}

For the coefficient algebra $A$, the kernel space is precisely
\[
 K(A,D)=\prod_\alpha u^\alpha D=B^1.
\]
Each $D$-valued coefficient is discrete, as in the continuous cochain
construction. Let $T:H\to A$ have coefficients
$T(e_i)=\sum_\beta t_{\beta i}u^\beta$.
For each $\beta$, only finitely many $t_{\beta i}$ are nonzero.
Define
\begin{equation}
 \delta_TF=\sum_iT(e_i)\partial_{u_i}F.
 \label{eq:op-derivation}
\end{equation}
For an output coefficient $\alpha$, this means the finite sum
\[
 (\delta_TF)_\alpha
  =\sum_{\beta+\gamma=\alpha}\sum_i
       t_{\beta i}(\gamma_i+1)F_{\gamma+\mathbf e_i}.
\]
It proves that $\delta_T$ is defined and continuous. The product rule
holds on coordinate polynomials, and continuity extends it to $A$.
Conversely, a continuous derivation restricts to a continuous map $H\to A$.
Formula~\eqref{eq:op-derivation} agrees with it on polynomials, hence
everywhere by density. Thus
\[
 \mathcal L(H,A)\cong\operatorname{Der}_{\mathrm{cont}}(A).
\]
This is a topological isomorphism when derivations inherit the operator
topology. Fixing the finitely many rows $T_\beta$ with $\beta\leq\alpha$ componentwise
fixes the entire $\alpha$-th row of $\delta_T$. Restriction is continuous
in the opposite direction.

An endomorphism of $H$ gives a derivation with coefficients linear in $u$.
On this subspace,
\[
 [\delta_S,\delta_T]=\delta_{ST-TS},
\]
because both sides agree on the linear functions $O_f$.
The identity kernel~\eqref{eq:op-identity} becomes
\[
 \mathcal E=\sum_i u_ic^i,\qquad
 \mathcal E(u^\alpha)=w(\alpha)u^\alpha,
 \qquad w(\alpha)=\sum_i\alpha_i.
\]
It is the Euler derivation. It acts as the identity on $H\subset A$,
while $\mathcal E(1)=0$.

Its operator degree on $H$ is zero. As a cochain it has coefficient
weight one and degree one. Before the cotangent regrading, $u_i$ has
degree two and $c^i$ has degree one, so $\mathcal E$ has degree three.
The contraction bracket has degree $-3$ and weight $-1$ before
regrading, and degree $-1$ afterwards.

Equation~\eqref{eq:op-diagonal} fixes the identity and hence fixes
$\mathcal E$. This invariance is different from cochain closedness.
The Chevalley--Eilenberg differential gives
\[
 (d\mathcal E)(f,g)
   =f\cdot O_g-g\cdot O_f-O_{[f,g]}=O_{[f,g]}.
\]
With the bivector $\Pi$ constructed in
Chapter~\ref{chap:hamiltonian-symmetry-cochains}, this is $d\mathcal E=\Pi$.
An arbitrary map $H\to A$ has no associative composition with another
such map. The derivation realization supplies a commutator instead.
Even the composition $\mathcal E^2$ is not a derivation: on weight $r$
it acts by $r^2$.

\section{Finite-rank trace and a disappearing boundary}

Let $\mathcal F(H)$ consist of continuous endomorphisms with
finite-dimensional image. Then
\[
 \mathcal F(H)=H\otimes D.
\]
To prove the nontrivial inclusion, let $W=T(H)$ be finite-dimensional.
The descending spaces $W\cap U_N$ eventually vanish: their dimensions
stabilize, and their intersection is zero. Some projection $W\to H_N$
is therefore injective. It follows that
$\ker T=\ker(q_NT)$ is open. In a basis of $W$, the coordinate
functionals of $T$ belong to $D$, giving its finite tensor expression.

The trace is
\begin{equation}
 \operatorname{tr}\!\left(\sum_{j=1}^r v_j\otimes\lambda_j\right)
       =\sum_{j=1}^r\lambda_j(v_j).
 \label{eq:op-finite-trace}
\end{equation}
It is independent of the expression. Restrict the operator to any
finite-dimensional subspace containing its image. The usual matrix trace
gives~\eqref{eq:op-finite-trace}; enlarging the subspace adds a zero
diagonal block. Finite-rank operators form a two-sided ideal, and the
rank-one formula proves
\[
 \operatorname{tr}(TF)=\operatorname{tr}(FT)
 \qquad(F\in\mathcal F(H),\ T\in\operatorname{End}_{\mathrm{cont}}(H)).
\]

The trace is not continuous in the inherited topology. Indeed,
\[
 P_n=e_n\otimes\epsilon^n\longrightarrow0,
 \qquad\operatorname{tr}(P_n)=1.
\]
Every fixed output row eventually vanishes, but the scalar does not
converge to zero in the discrete field. Therefore no continuous linear
extension of this trace to the completion exists. The same obstruction
is expressed by the failed uniformity condition for the evaluation
pairing in~\eqref{eq:op-uniform}.

On each stage $H\otimes H_M^\vee$, the trace uses only finitely
many coordinates of finitely many vectors in $H$, so it is continuous.
It is therefore continuous for the final linear topology of the union.
Its failure above proves that this topology differs from the tensor
topology whose completion represents all continuous operators.

Cyclicity fails for a stronger algebraic reason. Define the left and
right shifts by
\[
 L(e_0)=0,\quad L(e_n)=e_{n-1}\ (n\geq1),\qquad R(e_n)=e_{n+1}.
\]
Their row-finite matrices give continuous operators, and
\begin{equation}
 LR=1,\qquad RL=1-P_0,\qquad[L,R]=P_0.
 \label{eq:op-shift}
\end{equation}
A cyclic linear functional has $\tau(ST)=\tau(TS)$ and hence
$\tau(P_0)=0$. Extension of the finite-rank trace would give one.

On the finite space spanned by $e_0,\ldots,e_{m-1}$, the truncated
shifts instead satisfy
\[
 [L_m,R_m]=P_0-P_{m-1}.
\]
Their trace is zero. Passing to the infinite operator topology removes
$P_{m-1}$, whose trace remains one. The boundary term explains why
the finite-dimensional trace identity cannot pass to this completion.

\begin{proposition}
Every continuous endomorphism of $H$ is a commutator with the same
left shift $L$. Consequently every cyclic scalar linear functional
on $\operatorname{End}_{\mathrm{cont}}(H)$ is zero.
\end{proposition}
\begin{proof}
For a row-finite matrix $T$, define
\[
 X_{0j}=0,\qquad X_{i+1,j}=T_{ij}+X_{i,j-1},\qquad X_{i,-1}=0.
\]
Every row is finite by induction: it is a row of $T$ plus a shifted
preceding row of $X$. Thus $X$ is continuous. Matrix multiplication gives
$(LX)_{ij}=X_{i+1,j}$ and $(XL)_{ij}=X_{i,j-1}$.
The recurrence proves $[L,X]=T$, and cyclicity kills this commutator.
\end{proof}

A Hamiltonian translation can supply one side of the obstruction.
On the closed pure-$y$ subspace, use the basis
$y^{n+1}/(n+1)!$, $n\geq0$. Let $S$ shift this basis upward and
vanish on monomials containing $x$. The operator $a_x$ shifts the
displayed basis downward and preserves its complement. Hence
\[
 [a_x,S]=P_y,
\]
where $P_y$ projects onto $ky$. Any algebra containing these operators
has the same obstruction to a cyclic extension of finite-rank trace.

\section{A restricted algebra with an algebraic trace}

The trace obstruction depends on the operator algebra. A proper algebra
can retain every Hamiltonian vector field and still admit an algebraic
trace. Its topology remains a separate issue.

Put $R_0=k[[x,y]]$. Let $\mathcal D_0$ consist of finite-order
formal differential operators
\[
 P=\sum_{|\alpha|\leq m}a_\alpha(x,y)\partial^\alpha,
 \qquad P(1)\in k.
\]
Here $\alpha=(a,b)$ and $\partial^\alpha=\partial_x^a\partial_y^b$.
The last condition says that the zeroth-order coefficient is constant.
Such operators preserve constants and descend to $H$. They are continuous
in the formal topology. They form a unital algebra, since
$P(1)=c$, $Q(1)=d$ imply $PQ(1)=cd$.
Every Hamiltonian vector field $-f_y\partial_x+f_x\partial_y$ belongs
to this algebra.

\begin{lemma}
If a finite-order formal differential operator $Q$ annihilates
$\mathfrak m^N$, then $Q=0$.
\end{lemma}
\begin{proof}
For order zero, its coefficient $b$ satisfies $bx^N=0$ and vanishes
because $R_0$ is an integral domain. For positive order, $[Q,x]$ and
$[Q,y]$ annihilate the same ideal and have smaller order. Induction applies.
In normal form,
\[
 [Q,x]=\sum_\alpha\alpha_xb_\alpha
                   \partial^{\alpha-\mathbf e_x},
\]
and the analogous formula holds for $y$. Characteristic zero forces
every positive-order coefficient to vanish. The order-zero argument
then finishes the proof. This induction also proves uniqueness of a
normal-form expression as an operator.
\end{proof}

Suppose the descent of $P\in\mathcal D_0$ has finite rank. Its open
kernel gives $P(\mathfrak m^N)\subset k$ for some $N$. Thus
$\partial_xP$ annihilates $\mathfrak m^N$ and is zero by the lemma.
This forces $P=0$. Indeed, if $P$ has order $m$, form its highest-order
polynomial
\[
 \sigma_m(P)=\sum_{|\alpha|=m}a_\alpha\xi^\alpha
             \in R_0[\xi_x,\xi_y].
\]
The order-$(m+1)$ polynomial of $\partial_xP$ is
$\xi_x\sigma_m(P)$, nonzero in this integral domain.
Consequently descent is faithful and
$\mathcal D_0\cap\mathcal F(H)=0$.

The vector-space direct sum
\[
 \mathcal R=\mathcal D_0\oplus\mathcal F(H)
         \subset\operatorname{End}_{\mathrm{cont}}(H)
\]
is an associative algebra, because $\mathcal F(H)$ is an ideal.
The map $\chi(P)=P(1)$ is a scalar algebra homomorphism. For each $c\in k$,
define
\begin{equation}
 \tau_c(P+F)=c\chi(P)+\operatorname{tr}(F).
 \label{eq:op-restricted-trace}
\end{equation}
The direct decomposition makes this well-defined. It is cyclic:
$\chi([P,Q])=0$, and every mixed or finite-rank commutator has trace
zero by~\eqref{eq:op-finite-trace}. It extends finite-rank trace and
assigns $\tau_c(1)=c$.

This algebra contains all Hamiltonian vector fields but excludes the
operator $S$ above. Otherwise cyclicity would imply
$0=\tau_c([a_x,S])=\tau_c(P_y)=1$.
It is dense in the full operator algebra because it contains the dense
finite-rank ideal. Its trace is discontinuous, as the sequence $P_n$
still shows, and it has no cyclic extension to the completion.
Every positive-length product of Hamiltonian vector fields annihilates
the constant function. Formula~\eqref{eq:op-restricted-trace} therefore
assigns zero to all such products.

\section{Continuous invariant coefficients and quantum matter}

There is an additional obstruction even before choosing a trace.
Call a scalar multilinear form $\Phi:H^r\to k$ invariant if
\[
 \sum_{j=1}^r
 \Phi(f_1,\ldots,[a,f_j],\ldots,f_r)=0
 \qquad(a,f_1,\ldots,f_r\in H).
\]

\begin{proposition}
Every jointly continuous invariant scalar multilinear form on $H^r$
vanishes when $r\geq1$. Symmetry of the form is not required.
\end{proposition}
\begin{proof}
A continuous scalar multilinear form factors through finite coefficient
quotients in every argument. To see this directly, proceed by induction
on arity. At arity one, the kernel is open. At larger arity, continuity
at zero gives open subspaces whose product is annihilated. Choose
finite-dimensional complements. Fixing a basis vector in one complement
leaves a continuous form of smaller arity. Induction gives finite
coefficient bounds for each of these finitely many forms. Intersect their
bounds with the original open subspaces, and expand every argument into
its open and complementary parts. Each sufficiently small subspace then
annihilates the form in its individual slot.

The first-slot annihilator
\[
 J=\{f:\Phi(f,f_2,\ldots,f_r)=0\text{ for all }f_2,\ldots,f_r\}
\]
is therefore open. Invariance makes it a Lie ideal. If it contains
$\mathfrak m^N$, the formal antiderivative
\[
 I_y^Nf=\sum_{a+b\geq1}\frac{b!}{(b+N)!}f_{ab}x^ay^{b+N}
\]
belongs to $J$ and satisfies $(\operatorname{ad}_x)^NI_y^Nf=f$.
The ideal property gives $f\in J$ for every $f$, so $J=H$ and $\Phi=0$.
\end{proof}

Suppose a representation $\rho$ and a cyclic trace define
\[
 \Phi(f_1,\ldots,f_r)
       =\tau\bigl(\rho(f_1)\cdots\rho(f_r)\bigr).
\]
Expanding the trace of the commutator with $\rho(a)$ proves invariance.
If this scalar coefficient is jointly continuous for the formal topology,
the proposition forces it to vanish. In particular, these assumptions
exclude a nonzero continuous cubic trace polynomial.

The completed identity kernel supplies an actual operator realization.
An infinite-dimensional quantum current theory still needs a matter
space and paired space, contraction tensors, and a trace or regularization
on all the operator products that occur. Their continuity and convergence
must hold together with the spatial propagator integrals. The adjoint
choice does not reduce those products to finite rank:
$(a_x)^r I_y^r f=f$ makes every positive power of $a_x$ surjective.
The restricted trace above supplies
an algebraic choice, while its discontinuity prevents it from closing
completed kernels by a continuous scalar contraction. A quantum
construction must resolve these precise trace and contraction conditions.

Positive representations require an involution and a positive functional
in addition to a cyclic scalar contraction.
For finite matrix kernels, the positive trace gives a Hilbert space
completion but has no positive extension of finite mass to the unitization.
Chapter~\ref{chap:weil-operator-domains} proves this comparison and
formulates the arithmetic domain question using the zeta explicit formula.
